\documentclass[a4paper, 11pt]{article}

\usepackage[swedish]{babel}
\usepackage[utf8]{inputenc}
\usepackage[T1]{fontenc}
\usepackage{lmodern}

\usepackage{tikz}

\tikzset{
  excl/.style = {circle, draw=black, inner sep=0pt, minimum size=4pt}
}

\tikzset{
  incl/.style = {circle, draw=black, fill=black, inner sep=0pt, minimum size=4pt}
}


\usepackage{pgfplots}


\usepackage{graphicx}
\usepackage{caption}
\usepackage{subcaption}

\usepackage{enumitem}

\usepackage{multicol}

\usepackage{amsmath}
\usepackage{amsfonts}

\begin{document}

\title{
    \textbf{MM2001: Inläming 3}
}
\author{Johan Montelius}
\date{HT 2026}

\maketitle

\section*{i)}

Låt g vara funktionen $g(x) = x e^{−3/x}$ där $x \neq 0$

\begin{itemize}
  \begin{enumerate}[label=(\alph*)]
    
  \item Ange alla eventuella nollställen, lokala extrempunkter,
    asymptoter samt undersök funktionens beteende på randen av
    definitionsmängden.  Ange även i vilka intervall funktionen är
    växande, respektive avtagande.

  \item Gör en kvalitativ skiss av grafen till funktionen. Motivera
    noggrant hur du har kommit fram till din skiss!
\end{enumerate}

Funtion $g(x)$ är lika ned noll omm: $x$ är lika med noll eller då
$e^{-3/x}$ är lika med noll. Funktionen är dock inte definerad för $x = 0$ (eftersom $-3/0$ är odefinerat och uttrycket $e^t$ är alltid positivt. Detta ger att funktionen inte har några nollställen. 

Tar vi derivatan av funktionen $g(x)$ får vi: $$ g'(x) = e^{-3/x} + x -3

\subsection*{(a)}

\begin{itemize}
\item Nollställen: funktionen har inga nollställen.
\item Extrempunkter: $g'(x) = 
\end{itemize}




\subsection*{(b)}



\section*{(ii)}


- Vad menas med att en funktion $h : D \rightarrow R$ är växande? (Dvs
frågan är alltså efter definitionen av begreppet växande)

- Ange alla satser från boken PB1 som du har använt för att bestämma i
vilka intervall funktionen g ovan är växande. Beskriv kort hur har du
använt satserna!


\section*{(iii)}

Lägg särskilt fokus på funktionens beteende i en omgivning av punkten
$x = 0$. Har funktionen en (ensidig) derivata i $x = 0$ om vi sätter $g(0)
= 0$? Hur återspeglas ditt svar i skissen av grafen?


\end{document}


